\begin{abcliste}
\abc Die Tabellen zur Sonne (FoTaBe S.~211), zur Erde (S.~209) und zum Mond (S.~210) geben
\begin{align}
 m_\text{Sonne} &= \mSo \approx \result{\mSorP{0}{3}} \\
 m_\text{Erde} &= \mEa \approx \result{\mEarP{0}{3}} \\
 m_\text{Mond} &= \mMo \approx \result{\mMorP{0}{3}}
\end{align}

\abc Die Tabelle der Elemente (FoTaBe S.~228\,f.) gibt die molaren Massen in g/mol. Ihre Zahlenwerte sind die Atommassen in atomaren Masseneinheiten: $A_\text{He} = \AHeO$ und $A_\text{Pb} = \APbO$, mit $u = \uMaO$ (FoTaBe S.~218). Damit ist
\begin{align}
 m_\text{He} &= \mHeF = \AHe\cdot\uMa = \mHe \approx \result{\mHeP{0}{3}} \\
 m_\text{Pb} &= \mPbF = \APb\cdot\uMa = \mPb \approx \result{\mPbP{0}{3}}
\end{align}

\abc Die Ruhemasse des Elektrons (FoTaBe S.~218) ist
\begin{align}
 m_\text{e} &= \mel \approx \result{\melrP{0}{3}}
\end{align}
Für das Myon gibt die Tabelle der instabilen Teilchen (FoTaBe S.~196) die Ruheenergie $E_\mu = \EmuO$ an. Mit $m_\mu = E_\mu/c^2$ und dem Umrechnungsfaktor $\kMeVO$ aus der Tabelle zur Äquivalenz von Masse- und Energieeinheiten (FoTaBe S.~186) ist
\begin{align}
 m_\mu &= \mmuF = \Emu\cdot\kMeV = \mmu \approx \result{\mmuP{0}{3}}
\end{align}
\end{abcliste}
